\begin{figure}[t]
  \begin{panel}{(a)}{\textwidth}
    \setlength\plotwidth{0.6\textwidth}%
    \setlength\plotheight{36mm}%
    \small
    \centerline{\tikzsetnextfilename{good_fit_trace_posterior}%
\begin{tikzpicture}
  \setlength\plotwidth{0.7\linewidth}
  \argsinput[
    tracecsv=figures/data/posteriors/r1_10_traces_and_fits_revision.csv,
    traceintensitycol=N4_trace_good,
    zcol=N4_fit_good,
    plotlabel={$\noexpand\n=4$}]{figures/plots/fancy_trace.plot}
  \setlength\plotwidth{0.2\plotwidth}
  \pgfplotsset{every axis/.style={
    anchor=south west,
    at={($(trace.south east)+(6mm,0)$)},
    xtick distance=1
  }}
  \argsinput[
    posteriorcsv=figures/data/posteriors/r1_10_good_fit_posterior.csv,
    posteriorcol=n4_prob,
    xmin=1,
    xmax=8
  ]{figures/plots/p_n_given_x.plot}
  \node[above=0mm of posterior] {posterior};
\end{tikzpicture}
}
  \end{panel}
  \begin{panel}{(b)}{\textwidth}
    \setlength\plotwidth{0.6\textwidth}%
    \setlength\plotheight{36mm}%
    \small
    \centerline{\input{figures/panels/bad_fit_posterior.tikz.tex}}
  \end{panel}
  \caption{%
    \panelref{a} Experimental trace with a true count of 4, and the blinx 
    estimated posterior, identifying 4 as the most likely count. The posterior 
    displays a relatively high likelihood that the count could be $>4$, which is likely 
    due to increased noise.
    \panelref{b} Another experimental trace with a true count of 4, where the blinx 
    estimated posterior incorrectly finds that 3 is the most likely count. Posteriors 
    that are confidently incorrect such as this one reduce the average probability of 
    the correct count in \figref{fig:results:experimental}e.
  }
  \label{fig:exp_posterior}
\end{figure}
